\section{Eficiencia}
\label{sec:eficiencia}

\subsection{Tamaño y memoria}

Las cifras que no dependen de cómo se mida --- número de parámetros, tamaño
en disco, número de capas --- se resumen en la Tabla~\ref{tab:estatico}. Son
las que justifican por qué se esperaría que DistilBERT fuese más barato; las
dos subsecciones siguientes comprueban si esa expectativa se cumple en
tiempo de ejecución.

\begin{table}[htbp]
\centering
\small
\begin{tabular}{lrrr}
\toprule
 & \textbf{BERT-base} & \textbf{DistilBERT-base} & \textbf{Diferencia} \\
\midrule
Parámetros                    & \num{109483778} & \num{66955010} & $-\SI{38,8}{\percent}$ \\
Tamaño en fp32                & \SI{417,7}{\mega\byte} & \SI{255,4}{\mega\byte} & $-\SI{38,8}{\percent}$ \\
Capas \textit{transformer}    & 12 & 6 & $-\SI{50}{\percent}$ \\
Memoria GPU en inferencia     & $\approx$\SI{1720}{\mega\byte} & $\approx$\SI{1060}{\mega\byte} & $-\SI{38}{\percent}$ \\
Tiempo de entrenamiento (3 datasets) & \SI{29,3}{min} & \SI{16,4}{min} & $-\SI{44}{\percent}$ \\
\bottomrule
\end{tabular}
\caption{Métricas estáticas de eficiencia. La reducción de parámetros
(\SI{38,8}{\percent}) es menor que la de capas (\SI{50}{\percent}) porque la
matriz de \textit{embeddings}, que no se reduce, representa una fracción
importante del total.}
\label{tab:estatico}
\end{table}

\subsection{Por qué la latencia hubo que medirla aparte}

Cada corrida de \textit{fine-tuning} guarda un campo
\texttt{latencia\_ms\_media}, medido al terminar ese entrenamiento y con la
longitud de secuencia de \emph{su} dataset. Esos números \textbf{no son
comparables entre sí}, y al ponerlos juntos producen un resultado absurdo:
BERT en SST-2 marca \SI{3,16}{\milli\second} y DistilBERT en SST-2 marca
\SI{3,50}{\milli\second}. El modelo con el doble de capas saldría ``más
rápido''.

Detectar esto y no reportarlo sin más es parte del resultado. La causa, como
se ve enseguida, es que con \textit{batch} $=1$ la GPU no está calculando sino
esperando a que se lancen los \textit{kernels}, de modo que la medida no
refleja la arquitectura sino el ruido del lanzamiento. La solución fue un
módulo aparte (\texttt{src/benchmark.py}) que mide a los dos modelos en la
\textbf{misma rejilla} de (\textit{batch}, longitud), en la misma GPU y en el
mismo proceso. La latencia no depende de los pesos ajustados sino de la
arquitectura, así que se mide sobre los \textit{checkpoints} pre-entrenados
con entradas sintéticas.

\subsection{Latencia en condiciones controladas}

\begin{table}[h]
\centering
\footnotesize
\setlength{\tabcolsep}{4.5pt}
\begin{tabular}{rr rr c rr c}
\toprule
\textbf{Batch} & \textbf{Long.} & \textbf{BERT} & \textbf{DistilBERT} &
\textbf{Acel.} & \textbf{BERT} & \textbf{DistilBERT} & \textbf{Ratio} \\
 & & (ms) & (ms) & & (MB) & (MB) & mem. \\
\midrule
1  & 64  &   5,56 &  2,60 & \best{2,14$\times$} & 429,9 & 268,6 & 0,62$\times$ \\
1  & 128 &   5,61 &  2,61 & 2,15$\times$ & 433,0 & 271,6 & 0,63$\times$ \\
1  & 256 &   5,11 &  2,74 & 1,86$\times$ & 436,1 & 274,8 & 0,63$\times$ \\
8  & 64  &   7,72 &  4,01 & 1,93$\times$ & 445,9 & 284,5 & 0,64$\times$ \\
8  & 128 &  13,62 &  6,95 & 1,96$\times$ & 461,5 & 299,5 & 0,65$\times$ \\
8  & 256 &  26,59 & 13,39 & 1,99$\times$ & 493,9 & 332,5 & 0,67$\times$ \\
32 & 64  &  24,97 & 12,62 & 1,98$\times$ & 493,9 & 332,5 & 0,67$\times$ \\
32 & 128 &  49,85 & 25,10 & 1,99$\times$ & 559,9 & 398,6 & 0,71$\times$ \\
32 & 256 &  98,10 & 49,25 & 1,99$\times$ & 692,0 & 530,6 & 0,77$\times$ \\
\bottomrule
\end{tabular}
\caption{Latencia y pico de memoria en la misma rejilla (RTX A6000, fp32,
mismo proceso). DistilBERT es \num{2}$\times$ más rápido en los nueve puntos.}
\label{tab:benchmark}
\end{table}

\begin{figure}[h]
\centering
\includegraphics[width=\textwidth]{benchmark_latencia.pdf}
\caption{Latencia frente a longitud de secuencia, un panel por tamaño de
\textit{batch}. Se usan paneles separados en lugar de un solo gráfico porque
las latencias de \textit{batch} 1 y \textit{batch} 32 difieren en un orden de
magnitud y compartir eje aplastaría la curva inferior.}
\label{fig:benchmark}
\end{figure}

Medido así, DistilBERT es \best{2$\times$} más rápido en los nueve puntos de
la rejilla, que es exactamente lo que predice pasar de 12 capas a 6. En
rendimiento, con \textit{batch} 32 y 64 \textit{tokens}: \num{2535} frente a
\num{1282} muestras por segundo.

\paragraph{Los dos regímenes.} La Figura~\ref{fig:benchmark} muestra además
por qué los números por corrida no servían: \textbf{con \textit{batch} $=1$ la
latencia es plana} respecto de la longitud de secuencia: BERT marca
\num{5,56}, \num{5,61} y \num{5,11}~ms para 64, 128 y 256
\textit{tokens}, y DistilBERT \num{2,60}, \num{2,61} y \num{2,74}. Multiplicar por cuatro la longitud no
cuesta nada porque la GPU no está calculando: está esperando a que se lancen
los \textit{kernels}. Solo a partir de \textit{batch} 8 la latencia empieza a
seguir a la longitud, que es cuando la medida significa algo.

La conclusión práctica es que la ventaja de DistilBERT se cobra en los dos
regímenes, pero por motivos distintos: menos \textit{kernels} que lanzar
cuando se sirve petición a petición, y la mitad de cálculo cuando se procesa
por lotes.

\paragraph{La memoria no escala igual que los parámetros.} El ratio de
memoria pasa de \num{0,62}$\times$ con \textit{batch} 1 a \num{0,77}$\times$
con \textit{batch} 32 y 256 \textit{tokens}. Con \textit{batch} pequeño la
memoria está dominada por los pesos, que sí siguen la proporción de
parámetros; al crecer el \textit{batch}, las activaciones pasan a dominar, y
esas dependen de la longitud de secuencia y del tamaño del lote --- iguales
para los dos modelos --- más que del número de capas. En despliegues con
lotes grandes, la ventaja de memoria de DistilBERT se estrecha.

\subsection{Las tres dimensiones a la vez}

\begin{figure}[h]
\centering
\includegraphics[width=0.78\textwidth]{burbujas_params_acc.pdf}
\caption{Parámetros frente a \textit{accuracy} de test; el área de cada
burbuja es la latencia por muestra. Resume el intercambio completo: cuánto
ocupa, cuánto acierta y cuánto tarda cada combinación de modelo y tarea.}
\label{fig:burbujas}
\end{figure}

La Figura~\ref{fig:burbujas} condensa el argumento del informe. Los puntos de
DistilBERT están a la izquierda (menos parámetros) y en la misma banda
vertical que los de BERT (\textit{accuracy} comparable), con burbujas de área
similar o menor. La distancia vertical entre las dos columnas es pequeña en
todos los casos y prácticamente nula en AG News.
